% portada_algebra_lineal2.tex ─────────────────────────────────────────────────
% Portada Álgebra Lineal 2 — Universidad de Guanajuato
%
% ─── EDITAR AQUÍ ─────────────────────────────────────────────────────────────
\newcommand{\numtarea}{1}
\newcommand{\nomprofesor}{\rule{5cm}{0.3pt}}
% ─────────────────────────────────────────────────────────────────────────────

\documentclass[tikz, border=0pt]{standalone}

\usepackage[T1]{fontenc}
\usepackage[utf8]{inputenc}
\usepackage[spanish]{babel}
\usepackage{ebgaramond}
\usepackage[default]{sourcesanspro}
\usepackage{xcolor}
\usepackage{pgfplots}
\pgfplotsset{compat=1.18}

\definecolor{ugblue}{HTML}{002D72}
\definecolor{uggold}{HTML}{B8952A}

\begin{document}
\begin{tikzpicture}

  \useasboundingbox (0,0) rectangle (210mm,297mm);
  \fill[white] (0,0) rectangle (210mm,297mm);

  %% ── Barras de color ─────────────────────────────────────────────────────────
  \fill[ugblue] (0,287mm) rectangle (210mm,297mm);
  \fill[uggold] (0,284mm) rectangle (210mm,287mm);
  \fill[uggold] (0,11mm)  rectangle (210mm,14mm);
  \fill[ugblue] (0,0)     rectangle (210mm,11mm);

  %% ── Encabezado institucional ────────────────────────────────────────────────
  \node[font=\fontsize{15.5}{19}\selectfont\sffamily\bfseries\color{ugblue},
        anchor=center] at (105mm,275mm)
    {\MakeUppercase{Universidad de Guanajuato}};
  \node[font=\fontsize{10.5}{14}\selectfont\rmfamily\itshape\color{gray!65},
        anchor=center] at (105mm,267.5mm)
    {Divisi\'on de Ciencias Naturales y Exactas};
  \node[font=\fontsize{10.5}{14}\selectfont\rmfamily\itshape\color{gray!65},
        anchor=center] at (105mm,261mm)
    {Departamento de Matem\'aticas};

  %% ── Divisor superior ────────────────────────────────────────────────────────
  \draw[ugblue!22, line width=0.4pt] (30mm,254.5mm) -- (101mm,254.5mm);
  \fill[uggold]
    (103mm,254.5mm)--(105mm,256.5mm)--(107mm,254.5mm)--(105mm,252.5mm)--cycle;
  \draw[ugblue!22, line width=0.4pt] (109mm,254.5mm) -- (180mm,254.5mm);

  %% ── Hiperboloide de una hoja  x²+y² − z² = R² ──────────────────────────────
  %% Parametrización: x = R·cosh(v)·cos(u),  y = R·cosh(v)·sin(u),  z = R·sinh(v)
  %% Representa la forma cuadrática con eigenvalores (+,+,−): LA2
  %% Colormap: azul (abierto inferior) → cian → verde (cintura) → amarillo → rojo
  \begin{axis}[
    at          = {(105mm,186mm)},
    anchor      = center,
    width       = 160mm,
    height      = 118mm,
    axis lines  = none,
    view        = {38}{26},
    z buffer    = sort,
    clip        = false,
    colormap    = {almap}{
      rgb255=(0,0,160)     % azul profundo  (extremo inferior)
      rgb255=(0,80,255)    % azul
      rgb255=(0,200,230)   % cian
      rgb255=(0,190,80)    % verde          (cintura, z≈0)
      rgb255=(200,210,0)   % amarillo-verde
      rgb255=(255,130,0)   % naranja
      rgb255=(200,0,20)    % rojo profundo  (extremo superior)
    },
    point meta min = -7,
    point meta max =  7,
  ]
    \addplot3[
      surf,
      shader      = faceted interp,
      samples     = 42,
      samples y   = 32,
      domain      = 0:360,          % u: ángulo azimutal (grados)
      domain y    = -2.1:2.1,       % v: parámetro de altura (radianes)
      line width  = 0.05pt,
      point meta  = z,
    ]
    % Hiperboloide: R = 1.6
    ({1.6*cosh(y)*cos(x)},
     {1.6*cosh(y)*sin(x)},
     {1.6*sinh(y)});
  \end{axis}

  %% ── Divisor inferior ────────────────────────────────────────────────────────
  \draw[ugblue!22, line width=0.4pt] (30mm,121mm) -- (101mm,121mm);
  \fill[uggold]
    (103mm,121mm)--(105mm,123mm)--(107mm,121mm)--(105mm,119mm)--cycle;
  \draw[ugblue!22, line width=0.4pt] (109mm,121mm) -- (180mm,121mm);

  %% ── Título ──────────────────────────────────────────────────────────────────
  \node[font=\fontsize{9}{11}\selectfont\sffamily\fontseries{l}\selectfont
        \color{gray!48}, anchor=center] at (105mm,114mm)
    {\MakeUppercase{Licenciatura en Matem\'aticas}};

  \node[font=\fontsize{42}{48}\selectfont\sffamily\bfseries\color{ugblue},
        anchor=center] at (105mm,97mm)
    {\'Algebra Lineal~2};

  \node[font=\fontsize{27}{31}\selectfont\rmfamily\itshape\color{uggold},
        anchor=center] at (105mm,81mm)
    {Tarea~\numtarea};

  %% ── Separador y datos ───────────────────────────────────────────────────────
  \draw[ugblue!13, line width=0.35pt] (30mm,70.5mm) -- (180mm,70.5mm);

  \node[anchor=west,
    font=\fontsize{7.5}{9}\selectfont\sffamily\color{gray!52}]
    at (30mm,63mm) {\MakeUppercase{Alumno}};
  \node[anchor=west,
    font=\fontsize{7.5}{9}\selectfont\sffamily\color{gray!52}]
    at (120mm,63mm) {\MakeUppercase{NUA}};
  \node[anchor=west,
    font=\fontsize{15}{18}\selectfont\rmfamily\color{black!85}]
    at (30mm,55.5mm) {Ricardo Le\'on Mart\'inez};
  \node[anchor=west,
    font=\fontsize{15}{18}\selectfont\rmfamily\color{black!85}]
    at (120mm,55.5mm) {827328};

  \draw[ugblue!10, line width=0.3pt] (30mm,48mm) -- (180mm,48mm);

  \node[anchor=center,
    font=\fontsize{7.5}{9}\selectfont\sffamily\color{gray!52}]
    at (105mm,42mm) {\MakeUppercase{Profesor(a)}};
  \node[anchor=center,
    font=\fontsize{15}{18}\selectfont\rmfamily\color{black!85}]
    at (105mm,34.5mm) {\nomprofesor};

  \node[font=\fontsize{7.5}{9}\selectfont\sffamily\color{gray!36},
        anchor=center] at (105mm,21mm)
    {\MakeUppercase{Universidad de Guanajuato~·~Campus Guanajuato}};

\end{tikzpicture}
\end{document}
